\documentclass[a4paper,12pt]{article}

% Настройки языка и шрифтов
\usepackage[T2A]{fontenc}
\usepackage[utf8]{inputenc}
\usepackage[russian]{babel}
\usepackage{amsmath,amsfonts,amssymb}
\usepackage{graphicx}
\usepackage{color}
\usepackage{hyperref}
\usepackage{geometry}
\usepackage{indentfirst}
\usepackage{longtable}
\usepackage{booktabs}
\usepackage{float}

\geometry{left=3cm,right=1.5cm,top=2cm,bottom=2cm}


\begin{document}
% Титульный лист
\begin{titlepage}
    \centering
    {\fontsize{14pt}{16pt}\selectfont МОСКОВСКИЙ АВИАЦИОННЫЙ ИНСТИТУТ\\(НАЦИОНАЛЬНЫЙ ИССЛЕДОВАТЕЛЬСКИЙ УНИВЕРСИТЕТ)}\\[0.5cm]
    {\fontsize{12pt}{14pt}\selectfont Институт №8 «Компьютерные науки и прикладная математика»}\\[4cm]
    
    {\fontsize{16pt}{18pt}\selectfont \textbf{Лабораторные работы}}\\
    {\fontsize{14pt}{16pt}\selectfont \textbf{по курсу «Информационный поиск»}}\\[5cm]
    
    \vfill
    \begin{flushright}
        \begin{minipage}{0.51\textwidth}
            Выполнил: Шелаев Серафим Ильич\\
            Группа: М8О-408Б-22\\
            Преподаватель: Кухтичев Антон Алексеевич
        \end{minipage}
    \end{flushright}
    
    \vfill
    {\large Москва, 2026}
\end{titlepage}

\tableofcontents
\newpage

%% ====================================================================
\section{Введение}
%% ====================================================================
В рамках курса «Информационный поиск» была разработана полнофункциональная поисковая система \textbf{SportSearch}, включающая подсистемы сбора данных, обработки текста (токенизация, стемминг), индексации, булева и ранжированного поиска. В качестве корпуса использованы статьи англоязычной Википедии по тематике «Спорт».

Ядро системы (индексация и поисковый движок) реализовано на языке C++ с использованием собственных структур данных:
\begin{itemize}
    \item \textbf{DynArray} — динамический массив (замена \texttt{std::vector}).
    \item \textbf{HashDict} — хеш-таблица с открытой адресацией и хеш-функцией.
\end{itemize}
Обвязка (скачивание статей, импорт в БД, веб-интерфейс) написана на Python. Хранение метаданных — MongoDB. Веб-интерфейс — Flask.

%% ====================================================================
\section{Лабораторная работа №1: Добыча корпуса документов}
%% ====================================================================

\subsection{Цель работы}
Собрать и проанализировать корпус текстовых документов объёмом не менее 30\,000 статей для последующей индексации и поиска.

\subsection{Описание источников}
В качестве предметной области выбран \textbf{спорт}. Источник данных — англоязычная \textbf{Википедия}. Статьи скачиваются через Wikipedia REST API по 40 тематическим запросам:
\texttt{football}, \texttt{basketball}, \texttt{tennis}, \texttt{Olympic Games}, \texttt{FIFA World Cup}, \texttt{ice hockey}, \texttt{swimming sport}, \texttt{boxing}, \texttt{Formula One}, \texttt{rugby union}, \texttt{cycling sport}, \texttt{golf sport}, \texttt{athletics}, \texttt{gymnastics}, \texttt{wrestling sport}, \texttt{esports} и др.

Документы представляют собой энциклопедические статьи, содержащие заголовок и основной текст. Язык документов — английский.

\subsection{Структура документа}
«Сырой» документ (JSON от API) содержит:
\begin{itemize}
    \item Заголовок статьи (\texttt{title}).
    \item Полный текст в виде HTML (\texttt{extract}).
    \item URL страницы.
\end{itemize}

\subsection{Выделение текста}
Для выделения текста используется парсер, удаляющий HTML-теги, служебные блоки, таблицы и навигационные элементы. В итоговый текстовый документ включаются:
\begin{enumerate}
    \item Заголовок.
    \item Основной текст статьи (очищенный от разметки).
\end{enumerate}
Результат сохраняется в каталог \texttt{data/text/wikipedia/} — по одному \texttt{.txt}-файлу на статью. Исходные данные (JSON) сохраняются в \texttt{data/raw/wikipedia/}.

\subsection{Статистика корпуса}
\begin{table}[H]
\centering
\begin{tabular}{|l|c|}
\hline
\textbf{Параметр} & \textbf{Значение} \\
\hline
Количество документов & 31016 \\
\hline
Средний размер текста (на документ) & 12000 байт \\
\hline
Общий объём выделенного текста & 340 МБ \\
\hline
\end{tabular}
\caption{Статистика собранного корпуса}
\end{table}

\subsection{Проверка пригодности}
Поиск по Википедии доступен через встроенный поиск (\url{https://en.wikipedia.org/w/index.php?search=...}), а также через Google с ограничением \texttt{site:en.wikipedia.org}. Это позволяет использовать данные источники как альтернативы для сравнительной оценки качества поиска.

\textbf{Примеры запросов к существующим поисковикам:}
\begin{itemize}
    \item \texttt{FIFA World Cup history} — встроенный поиск Wikipedia находит основную статью, но не все связанные (отборочные турниры, история отдельных чемпионатов).
    \item \texttt{olympic gold medal records} — Google \texttt{site:en.wikipedia.org} выдаёт релевантные результаты, но встроенный поиск Википедии может путать с неспортивными значениями слова «record».
\end{itemize}

%% ====================================================================
\section{Лабораторная работа №2: Поисковый робот}
%% ====================================================================

\subsection{Цель работы}
Разработать поисковый робот (crawler) для автоматического сбора документов с поддержкой докачки и соблюдения вежливости.

\subsection{Архитектура и стек технологий}
Робот реализован на языке \textbf{Python}. Основные компоненты:
\begin{itemize}
    \item \textbf{YAML-конфигурация} (\texttt{conf.yaml}): параметры запуска (задержки, timeout, user-agent, список seed-запросов).
    \item \textbf{MongoDB}: хранение метаданных документов и очереди скачивания.
    \item \textbf{Requests}: HTTP-запросы с повторными попытками (retries).
    \item \textbf{Wikipedia API}: скачивание статей через \texttt{action=query} с параметром \texttt{extracts}.
\end{itemize}

\subsection{Алгоритм работы}
\begin{enumerate}
    \item \textbf{Инициализация}: чтение конфигурации, подключение к MongoDB.
    \item \textbf{Поиск статей}: для каждогоseed-запроса выполняется поиск через Wikipedia API (\texttt{list=search}). Найденные заголовки добавляются в очередь.
    \item \textbf{Загрузка}: для каждого заголовка запрашивается полный текст (\texttt{prop=extracts}). Между запросами выдерживается задержка (400 мс).
    \item \textbf{Дедупликация}: хеш SHA-256 используется для обнаружения дубликатов; повторно скачанные страницы обновляют метку \texttt{fetched\_at}.
    \item \textbf{Докачка}: при перезапуске робот продолжает с того места, где остановился, выбирая из базы ещё не скачанные URL.
\end{enumerate}

\subsection{Результат}
Робот успешно собрал корпус спортивных статей из Википедии. Реализована устойчивость к разрывам соединения и возможность остановки/продолжения без потери прогресса.

%% ====================================================================
\section{Лабораторная работа №3: Токенизация}
%% ====================================================================

\subsection{Цель работы}
Реализовать разбиение текста на токены (слова) для построения поискового индекса.

\subsection{Правила токенизации}
Токенизатор реализован на C++ (файл \texttt{porter.hpp}, пространство имён \texttt{Porter}):
\begin{enumerate}
    \item \textbf{Разбиение}: текст сканируется посимвольно; токеном считается непрерывная последовательность алфавитно-цифровых символов (\texttt{isalnum}).
    \item \textbf{Нормализация}: все символы приводятся к нижнему регистру (\texttt{tolower}).
    \item \textbf{Стемминг}: к каждому токену применяется алгоритм Портера (см.~ЛР4).
\end{enumerate}

\subsection{Достоинства метода}
\begin{itemize}
    \item \textbf{Простота}: минимальные требования к памяти и вычислениям.
    \item \textbf{Высокая скорость}: однопроходный алгоритм $O(n)$.
    \item \textbf{Универсальность}: работает с любым текстом на латинице.
    \item \textbf{Унификация словоформ}: стемминг объединяет различные формы слова (players, playing, played $\to$ play).
\end{itemize}

\subsection{Недостатки метода}
\begin{itemize}
    \item \textbf{Потеря информации}: агрессивный стемминг может приводить к коллизиям (например, \texttt{university} $\to$ \texttt{univers}, при этом \texttt{universal} $\to$ \texttt{univers}).
    \item \textbf{Разбиение составных терминов}: дефисы являются разделителями, поэтому \texttt{100-meter} становится двумя токенами \texttt{100} и \texttt{meter}.
    \item \textbf{Числа}: даты и счёт матчей (\texttt{2-1}, \texttt{2024}) попадают в словарь как отдельные числовые токены.
\end{itemize}

\subsection{Примеры неудачной токенизации}
\begin{table}[H]
\centering
\begin{tabular}{|l|l|l|}
\hline
\textbf{Исходный текст} & \textbf{Результат} & \textbf{Проблема} \\
\hline
100-meter & 100, meter & Разбито на два токена \\
\hline
university & univers & Коллизия с «universal» \\
\hline
2--1 & 2, 1 & Счёт матча потерян \\
\hline
FIFA's & fifa, s & Апостроф разделяет слово \\
\hline
\end{tabular}
\caption{Примеры неудачной токенизации}
\end{table}

\textbf{Возможные улучшения:}
\begin{itemize}
    \item Добавить словарь составных терминов (n-grams) для спортивной лексики.
    \item Обрабатывать апострофы (\texttt{'s}, \texttt{'t}) как часть слова.
    \item Применить лемматизацию вместо стемминга для снижения коллизий.
\end{itemize}

\subsection{Статистика токенизации}
\begin{table}[H]
\centering
\begin{tabular}{|l|c|}
\hline
\textbf{Параметр} & \textbf{Значение} \\
\hline
Количество документов & 31016 \\
\hline
Общее количество токенов & 61094782 \\
\hline
Уникальных токенов (до стемминга) & 680753 \\
\hline
Среднее токенов на документ & 8227 \\
\hline
\end{tabular}
\caption{Статистика токенизации}
\end{table}

\subsection{Производительность}
\begin{table}[H]
\centering
\begin{tabular}{|l|c|}
\hline
\textbf{Параметр} & \textbf{Значение} \\
\hline
Время токенизации (1000 файлов) & 1,52 сек \\
\hline
Скорость токенизации & 8,32 КБ/сек \\
\hline
\end{tabular}
\caption{Производительность токенизатора}
\end{table}

\textbf{Зависимость от объёма}: время линейно зависит от размера входных данных ($T = O(n)$), так как алгоритм однопроходный. Возможные способы ускорения: многопоточная обработка, SIMD-инструкции, memory-mapped I/O.

%% ====================================================================
\section{Лабораторная работа №4: Стемминг}
%% ====================================================================

\subsection{Цель работы}
Реализовать морфологическую нормализацию (стемминг) для улучшения полноты поиска.

\subsection{Метод решения}
Реализован \textbf{алгоритм Портера} (Porter Stemmer) для английского языка на C++ (файл \texttt{porter.hpp}). Алгоритм последовательно применяет ряд правил для отсечения суффиксов и окончаний:
\begin{itemize}
    \item Шаг 1a: Обработка множественного числа (sses $\to$ ss, ies $\to$ i, s $\to$ $\varnothing$).
    \item Шаг 1b: Обработка окончаний ed, ing.
    \item Шаг 1c: Замена y на i.
    \item Шаг 2--4: Замена длинных суффиксов (ational $\to$ ate, izer $\to$ ize, и т.д.).
    \item Шаг 5: Удаление конечного e и двойных согласных (ll $\to$ l).
\end{itemize}

\subsection{Результаты}
Применение стемминга позволило сократить размер словаря примерно на 30--40\%, объединив различные словоформы в один терм.

\subsection{Интеграция в поисковую систему}
Стемминг применяется в двух местах:
\begin{enumerate}
    \item \textbf{При индексации}: каждый токен документа проходит через Porter Stemmer перед добавлением в индекс.
    \item \textbf{При поиске}: запрос пользователя обрабатывается тем же стеммером (\texttt{Porter::tokenize()}).
\end{enumerate}

Это обеспечивает \textbf{поиск без учёта словоформ}: запрос «players» найдёт документы, содержащие «player», «playing», «played» и другие формы.

\textbf{Примеры унификации:}
\begin{table}[H]
\centering
\begin{tabular}{|l|l|}
\hline
\textbf{Исходные слова} & \textbf{Результат стемминга} \\
\hline
championship, championships & championship \\
\hline
running, runner, runs, ran & run \\
\hline
players, playing, played, plays & play \\
\hline
swimming, swimmer, swam & swim \\
\hline
tournament, tournaments & tournament \\
\hline
\end{tabular}
\caption{Примеры унификации словоформ}
\end{table}

%% ====================================================================
\section{Лабораторная работа №5: Закон Ципфа}
%% ====================================================================

\subsection{Цель работы}
Проверить выполнение закона Ципфа на собранном корпусе документов.

\subsection{Методика}
Построен график распределения частот терминов в двойном логарифмическом масштабе (log-log plot):
\begin{itemize}
    \item По оси X: ранг слова (от самого частотного к редкому).
    \item По оси Y: частота встречаемости слова.
    \item Красная линия: теоретический закон Ципфа $f = \frac{C}{r}$, где $C$ --- частота самого частотного слова.
\end{itemize}

\subsection{Результаты}

\begin{figure}[H]
\centering
\includegraphics[width=0.9\textwidth]{zipf_sport.png}
\caption{Распределение частот терминов (закон Ципфа)}
\end{figure}

\textbf{Топ-10 самых частотных слов:}
\begin{table}[H]
\centering
\begin{tabular}{|c|l|r|}
\hline
\textbf{Ранг} & \textbf{Терм} & \textbf{Частота} \\
\hline
1 & the &  4493046 \\
2 & in   & 1733587 \\
3 & and & 1583866 \\
4 & of & 1275920 \\
5 & to & 1148356 \\
6 & a & 1110187 \\
7 & team & 817483 \\
8 & was & 539712 \\
9 & on & 536479 \\
10 & for & 498080 \\
\hline
\end{tabular}
\caption{Топ-10 частотных терминов}
\end{table}

Согласно закону Ципфа: $f \approx \frac{C}{r^s}$, где $s \approx 1$.

Полученный график демонстрирует линейную зависимость с наклоном, близким к $-1$, что подтверждает выполнение закона Ципфа для собранного корпуса спортивных статей.

\subsection{Причины расхождений}
На графике наблюдаются характерные отклонения от идеального закона:

\textbf{1. Область высоких рангов (частотные слова):}
\begin{itemize}
    \item Частота самых частотных слов (the, of, and) \textbf{ниже} предсказанной законом.
    \item Причина: стоп-слова распределены равномерно по всем документам, а не концентрированно.
\end{itemize}

\textbf{2. Область низких рангов (редкие слова):}
\begin{itemize}
    \item Наблюдается \textbf{ступенчатость} --- много слов с частотой 1, 2, 3.
    \item Частота \textbf{выше} предсказанной (толстый хвост).
    \item Причины: имена спортсменов, названия команд и стадионов, числа (годы, счёт), специфические термины (подвиды спорта).
\end{itemize}

\textbf{3. Средняя область:}
\begin{itemize}
    \item Хорошее соответствие закону Ципфа, типичное для естественных текстов.
\end{itemize}

\subsection{Выводы}
Закон Ципфа выполняется для собранного корпуса с типичными отклонениями. Отклонения объясняются особенностями спортивной тематики (обилие имён собственных и числовых данных).

%% ====================================================================
\section{Лабораторная работа №6: Булев индекс}
%% ====================================================================

\subsection{Цель работы}
Построить обратный индекс для булева поиска, используя собственные реализации структур данных и бинарный формат хранения.

\subsection{Внутреннее представление документа}
После токенизации документ представляется как последовательность термов (стеммированных токенов). Каждый терм — строка в нижнем регистре. Для каждого терма хранится позиция в документе (порядковый номер токена) — это обеспечивает координатный индекс для ЛР5.

\subsection{Реализация индексатора}
Индексатор написан на C++ (файл \texttt{idx\_main.cpp}) без использования STL контейнеров для хранения данных.

\subsubsection{Структуры данных}
\begin{itemize}
    \item \textbf{HashDict}: хеш-таблица с открытой адресацией и линейным пробированием, хеш-функция FNV-1a. Используется для накопления словаря терминов и постинг-листов в памяти.
    \item \textbf{DynArray}: динамический массив с автоматическим расширением (фактор роста $\times$2). Используется для хранения списков doc\_id и позиций.
\end{itemize}

\subsubsection{Метод сортировки}
Постинг-листы \textbf{не требуют явной сортировки}: документы обрабатываются последовательно (doc\_id = 0, 1, 2, ...), поэтому posting lists автоматически упорядочены по doc\_id.

\textbf{Достоинства:}
\begin{itemize}
    \item Нет накладных расходов на сортировку ($O(1)$ вместо $O(n \log n)$).
    \item Простота реализации.
\end{itemize}

\textbf{Недостатки:}
\begin{itemize}
    \item Требует последовательной обработки документов.
    \item При инкрементальном обновлении индекса потребуется merge-sort.
\end{itemize}

\subsubsection{Бинарный формат индекса (побайтовое описание)}
Индекс сохраняется в три файла:

\textbf{1. index.fwd (Прямой индекс):}
\begin{verbatim}
[0-3]   Magic: "FWRD" (4 байта)
[4-5]   Version: uint16 = 1 (2 байта)
[6-9]   DocCount: uint32 (4 байта)
[10..]  OffsetTable: uint64[DocCount] (8 байт × DocCount)
[...]   Данные документов:
        - url_len: uint16 (2 байта)
        - url: char[url_len]
        - title_len: uint16 (2 байта)
        - title: char[title_len]
        - doc_length: uint32 (4 байта) — длина в токенах
\end{verbatim}

\textbf{2. index.term (Словарь термов):}
\begin{verbatim}
[0-3]   Magic: "TERM" (4 байта)
[4-5]   Version: uint16 = 1 (2 байта)
[6-9]   TermCount: uint32 (4 байта)
[...]   Записи термов:
        - term_len: uint8 (1 байт)
        - term: char[term_len]
        - data_offset: uint64 (8 байт) — смещение в index.data
        - doc_freq: uint32 (4 байта) — количество документов
\end{verbatim}

\textbf{3. index.data (Обратный индекс / постинги):}
\begin{verbatim}
[0-3]   Magic: "DATA" (4 байта)
[4-5]   Version: uint16 = 1 (2 байта)
[...]   Постинг-листы для каждого терма:
        - doc_freq: VarByte
        - num_skips: VarByte
        - skip_table[num_skips]:
            - skip_doc_id: VarByte
            - skip_byte_offset: VarByte
        - entries[doc_freq]:
            - doc_id_delta: VarByte
            - term_freq: VarByte
            - positions[term_freq]: VarByte[] (delta-кодир.)
\end{verbatim}

\textbf{VarByte кодирование:} старший бит каждого байта — флаг продолжения (1 = ещё есть байты, 0 = последний байт). Остальные 7 бит — данные. Эффективно кодирует малые числа в 1 байт, большие — до 5 байт.

\textbf{Расширяемость:} поле Version позволяет изменять формат в будущих версиях с сохранением обратной совместимости.

\subsection{Результаты}
\begin{table}[H]
\centering
\begin{tabular}{|l|c|}
\hline
\textbf{Параметр} & \textbf{Значение} \\
\hline
Количество термов & 254605 \\
\hline
Средняя длина терма & 6.88 символов \\
\hline
Средняя длина токена (до стемминга) & 5.2 символов \\
\hline
Размер index.fwd & 6.98 МБ \\
\hline
Размер index.term & 21.74  МБ \\
\hline
Размер index.data & 28.7 МБ \\
\hline
Общий размер индекса & 57.42 МБ \\
\hline
\end{tabular}
\caption{Статистика индекса}
\end{table}

\textbf{Сравнение длины терма и токена:} Средняя длина терма (в словаре 6.88) обычно больше средней длины токена (по тексту 5.20), поскольку в тексте преобладают короткие высокочастотные слова (a, the, in), а в словаре каждый терм учитывается однократно. Стемминг также влияет на среднюю длину.

\subsection{Скорость индексации}
\begin{table}[H]
\centering
\begin{tabular}{|l|c|}
\hline
\textbf{Параметр} & \textbf{Значение} \\
\hline
Общее время индексации & 38 сек \\
\hline
Скорость (КБ/сек) & 6.4 \\
\hline
\end{tabular}
\caption{Скорость индексации}
\end{table}

\subsection{Анализ оптимальности}
\textbf{Текущие ограничения:}
\begin{itemize}
    \item Индекс строится целиком в памяти — ограничен объёмом RAM.
    \item Однопоточная обработка.
\end{itemize}

\textbf{Возможные оптимизации:} многопоточная токенизация, внешняя сортировка (external merge sort) для больших корпусов, memory-mapped I/O.

\textbf{Масштабируемость:}
\begin{itemize}
    \item \textbf{10$\times$ данных}: справится, потребуется $\approx$ 300 МБ RAM.
    \item \textbf{100$\times$ данных}: потребуется внешняя сортировка.
    \item \textbf{1000$\times$ данных}: распределённая система (MapReduce).
\end{itemize}

%% ====================================================================
\section{Лабораторная работа №7: Булев поиск}
%% ====================================================================

\subsection{Цель работы}
Реализовать выполнение булевых запросов (AND, OR, NOT) по индексу с поддержкой веб-интерфейса и CLI.

\subsection{Синтаксис запросов}
Парсер поддерживает:
\begin{itemize}
    \item \textbf{Пробел или \texttt{\&\&}} — логическое И (пересечение).
    \item \textbf{\texttt{||}} — логическое ИЛИ (объединение).
    \item \textbf{\texttt{!}} — логическое НЕ (исключение).
    \item \textbf{Скобки \texttt{( )}} — группировка операций.
\end{itemize}
Парсер реализован методом \textbf{рекурсивного спуска} и устойчив к переменному числу пробелов.

\subsection{Алгоритм поиска}
Поисковый движок (C++, файл \texttt{qry\_main.cpp}) работает следующим образом:
\begin{enumerate}
    \item \textbf{Загрузка}: Словарь загружается в \texttt{HashDict}, данные постингов (\texttt{index.data}) — целиком в память.
    \item \textbf{Токенизация запроса}: Выделение термов и операторов. Термы проходят стемминг через \texttt{Porter::stem()}.
    \item \textbf{Парсинг}: Построение дерева выражения методом рекурсивного спуска.
    \item \textbf{Выполнение}:
    \begin{itemize}
        \item \textbf{AND}: Пересечение сортированных списков (два указателя, $O(N+M)$).
        \item \textbf{OR}: Слияние списков ($O(N+M)$).
        \item \textbf{NOT}: Разность с множеством всех документов.
    \end{itemize}
\end{enumerate}

\subsection{Примеры запросов}
\begin{table}[H]
\centering
\begin{tabular}{|l|c|l|}
\hline
\textbf{Запрос} & \textbf{Результат} & \textbf{Описание} \\
\hline
\texttt{football} & 1516 док. & Простой поиск \\
\hline
\texttt{football \&\& world} & 987 док. & Пересечение (AND) \\
\hline
\texttt{basketball || volleyball} & 3178 док. & Объединение (OR) \\
\hline
\texttt{!football} & 29500 док. & Исключение (NOT) \\
\hline
\texttt{(football || rugby) \&\& world} & 1120 док. & Сложный запрос \\
\hline
\end{tabular}
\caption{Примеры булевых запросов}
\end{table}

\subsection{Веб-интерфейс}
Реализован веб-сервис на Flask (\texttt{server.py}):
\begin{itemize}
    \item \textbf{Главная страница} (\texttt{/}): Форма ввода поискового запроса.
    \item \textbf{Страница результатов} (\texttt{/search?q=...}): Форма поиска, количество найденных документов, время выполнения, до 50 результатов на страницу (заголовок + ссылка) и пагинация.
\end{itemize}

\subsection{CLI-утилита}
Реализована утилита командной строки (\texttt{search\_cli.py}), которая:
\begin{itemize}
    \item Запускает C++ движок (\texttt{bin/qry}) как фоновый процесс.
    \item Читает запросы из stdin (по одному на строку).
    \item Выводит результаты в stdout.
\end{itemize}

\subsection{Производительность}
\begin{table}[H]
\centering
\begin{tabular}{|l|c|}
\hline
\textbf{Метрика} & \textbf{Значение} \\
\hline
Загрузка индекса & 45 мс \\
\hline
Простой запрос (1 терм) & 1 мс \\
\hline
Сложный запрос (AND/OR/NOT) & 3 мс \\
\hline
\end{tabular}
\caption{Производительность поисковой системы}
\end{table}

\textbf{Примеры долгих запросов:}
\begin{itemize}
    \item Запросы с NOT на частотных термах (\texttt{!the}) — обход всех документов.
    \item Множественные OR (\texttt{a || b || c || d}) — большие объединения.
\end{itemize}

\subsection{Тестирование корректности}
Корректность проверялась следующими методами:
\begin{enumerate}
    \item \textbf{Инвариант NOT}: $|A| + |\neg A| = N$ (общее число документов).
    \item \textbf{Инвариант AND}: $|A \cap B| \leq \min(|A|, |B|)$.
    \item \textbf{Инвариант OR}: $|A \cup B| \geq \max(|A|, |B|)$.
    \item \textbf{Сравнение с grep}: Выборочная проверка результатов путём поиска терма в исходных текстах.
\end{enumerate}


%% ====================================================================
\section{Лабораторная работа №8: Ранжирование TF-IDF}
%% ====================================================================

\subsection{Цель работы}
Реализовать ранжированный поиск на основании схемы TF-IDF.

\subsection{Формулы}
Для каждой пары (терм $t$, документ $d$) вычисляется:
$$\text{TF}(t,d) = \frac{f_{t,d}}{|d|}$$
$$\text{IDF}(t) = \log\frac{N}{\text{df}(t)}$$
$$\text{Score}(d, q) = \sum_{t \in q} \text{TF}(t,d) \cdot \text{IDF}(t)$$

где $f_{t,d}$ --- частота терма $t$ в документе $d$, $|d|$ --- длина документа в токенах, $N$ --- общее число документов, $\text{df}(t)$ --- число документов, содержащих терм $t$.

\subsection{Интерпретация запросов}
\begin{itemize}
    \item Если запрос содержит \textbf{только термы через пробелы} — трактуется как \textbf{нечёткий} (free-text): результат ранжируется по TF-IDF, допускается неполное соответствие.
    \item Если запрос содержит \textbf{операторы} (\texttt{\&\&}, \texttt{||}, \texttt{!}, кавычки) — трактуется как \textbf{булев}: строгое соответствие, но порядок выдачи определяется TF-IDF.
\end{itemize}

\subsection{Примеры запросов}
\begin{table}[H]
\centering
\begin{tabular}{|l|l|l|}
\hline
\textbf{Запрос} & \textbf{Тип} & \textbf{Комментарий} \\
\hline
\texttt{olympic games history} & нечёткий & Ранжирование по TF-IDF \\
\hline
\texttt{football \&\& world} & булев & Строгое AND, порядок по TF-IDF \\
\hline
\texttt{best basketball players} & нечёткий & Допускается неполное совпадение \\
\hline
\texttt{"tour de france"} & фразовый & Точная фраза \\
\hline
\end{tabular}
\caption{Примеры запросов с ранжированием}
\end{table}

\textbf{Удачные запросы}: запросы из 2--3 специфических терминов (например, \texttt{wimbledon championship}) дают релевантные результаты на первых позициях, поскольку IDF высокочастотных слов мал, а специфических — велик.

\textbf{Неудачные запросы}: запросы из одних стоп-слов или очень коротких слов (например, \texttt{a the in}) дают практически случайный порядок выдачи, т.к. IDF таких слов близок к нулю.


%% ====================================================================
\section{Выводы}
%% ====================================================================
В ходе лабораторных работ была создана поисковая система \textbf{SportSearch} полного цикла:
\begin{enumerate}
    \item \textbf{Сбор данных}: краулер на Python скачивает статьи из Википедии по спортивной тематике.
    \item \textbf{Обработка текста}: токенизация и стемминг (алгоритм Портера) на C++.
    \item \textbf{Индексация}: координатный индекс с VarByte-сжатием и skip-указателями на C++.
    \item \textbf{Поиск}: булев (AND, OR, NOT), фразовый, proximity, TF-IDF ранжирование — на C++.
    \item \textbf{Интерфейс}: веб-интерфейс (Flask) и CLI-утилита.
\end{enumerate}

Реализация ядра на C++ без STL и применение VarByte-сжатия позволили достичь компактного индекса (57.42 МБ) и высокой скорости поиска (простой запрос $<$ 1 мс). Система успешно справляется с индексацией и поиском по корпусу из 31\,016 документов, содержащему более 61 миллиона токенов.

\end{document}
